\section{Introduction}

\plainMHCI{} molecules display intracellularly derived peptides at the surface of nucleated cells for surveillance by cytotoxic CD8\(^{+}\) T cells. This pathway enables immune recognition of infected and malignant cells and is central to epitope discovery, vaccine development, and T-cell-based cancer immunotherapy. In particular, the selection of peptides that are naturally presented by tumor cells is a prerequisite for personalized neoantigen vaccines and other antigen-directed therapies \cite{ott2017neoantigen,sahin2017mutanome,marcu2021atlas}. Large curated resources such as the Immune Epitope Database and tissue immunopeptidome databases further demonstrate the breadth of biomedical research that depends on reliable MHC ligand identification \cite{vita2025iedb,lemke2025pcidb}.

Computational prediction provides a practical first filter for the vast number of peptide--MHC combinations. Classical and pan-allelic methods such as NetMHCpan and ACME primarily estimate binding from peptide sequence and MHC information \cite{nielsen2007netmhcpan,hu2019acme}. With increasing mass-spectrometry eluted-ligand data, NetMHCpan-4.1, MHCflurry 2.0, and HLAthena extended this paradigm toward presentation by integrating binding, processing, and immunopeptidomic evidence \cite{reynisson2020netmhcpan,odonnell2020mhcflurry,sarkizova2020hlathena}. These methods established the transition from binding-centric prediction to broader presentation modeling.

Recent sequence-based work has diversified the representations used for peptide--MHC prediction. Anthem and MHCfovea support customized models or relate binding motifs to informative HLA residues \cite{mei2021anthem,lee2021mhcfovea}. TransPHLA, STMHCpan, and MHCSeqNet2 use attention-based architectures and learned HLA representations to extend coverage across alleles and peptide sequences \cite{chu2022transphla,ye2023stmhcpan,wongklaew2024mhcseqnet2}. These methods improve representation learning but mainly target peptide--MHC compatibility or general presentation.

A complementary trend incorporates richer sequence and three-dimensional information. CapsNet-MHC enriches sequence representation, while structure-conditioned networks, MHC-Fine, and geometric deep learning model peptide--MHC specificity from predicted or explicit complex structures \cite{kalemati2023capsnet,motmaen2023structure,glukhov2024mhcfine,marzella2024geometric}. Transfer and hierarchical progressive learning further address kinetic stability, immunogenicity, zero-shot binding, and antigen design \cite{fasoulis2024transfer,zhu2025hierarchical}. These approaches improve molecular interaction modeling and generalization, but tissue context generally remains outside their prediction target.

Presentation-oriented models increasingly learn from immunopeptidomic and biological context beyond binding affinity. RBM-MHC, MixMHCpred2.2, and BigMHC address ligand assignment, motif deconvolution, presentation, or immunogenicity ranking \cite{bravi2021rbmmhc,gfeller2023mixmhcpred,albert2023bigmhc}. ImmuneApp and HLApollo incorporate multi-allelic data, peptide flanks, molecular context, or processing signals, while NetMHCpan-4.2 extends transfer and structural modeling \cite{xu2024immuneapp,thrift2024hlapollo,nilsson2025netmhcpan}. These developments advance general presentation and epitope prioritization, but do not directly model tissue-specific preference; evaluation generally pools ligands, alleles, or samples rather than controlling peptide pairs within a tissue context.

Despite these advances, existing tools primarily predict peptide--MHC binding or general presentation rather than which peptide fragment is preferentially presented in a specified tissue. This distinction is biologically important because tissue-dependent cleavage, transport, trimming, and peptide loading can shape the fragments that reach the cell surface, and multi-tissue immunopeptidomic studies in humans and mice have identified marked tissue-associated repertoires, including different fragments derived from the same source proteins across tissues \cite{marcu2021atlas,schuster2018mouse,kubiniok2022constitutive}. Characterizing such preferences could support tissue-aware tumor-antigen prioritization and assessment of potential on-target, off-tumor recognition \cite{marcu2021atlas,lemke2025pcidb}. We therefore formulate tissue-specific preference prediction for MHC-I presentation as a relative question: when the source protein and MHC restriction are fixed, which of two peptides has stronger presentation evidence in a target tissue? TissuePMHC addresses this question using matched peptide pairs that also share the same cross-tissue occurrence count, shifting the prediction target from general presentation propensity to tissue-specific preference while complementing existing binding and presentation predictors.

To address this question, we develop TissuePMHC using matched human and mouse peptide data. Each comparison contains two peptides with the same MHC restriction, source protein, and overall frequency of observation across tissues. This design reduces obvious shortcuts and asks the model to identify which peptide has stronger presentation evidence in a particular tissue. TissuePMHC learns both broadly shared peptide patterns and patterns associated with individual MHC molecules. Our conclusions are limited to the tissue--MHC combinations represented in the available data.

The main contributions of this study are to define tissue-specific preference as a controlled peptide-comparison problem, construct matched benchmarks in two species, test whether the resulting signal can be explained by existing general presentation predictors, and examine the peptide features used by TissuePMHC. Together, these contributions provide a practical framework for studying how tissue context may influence which fragments of the same protein are presented by MHC-I.
